\section{Halting destroys every encoding}\label{sec:forgetting}

A classical controller can forget its initial state while a coupled
quantum system retains information elsewhere. We prove the required
quantum statement through initial dependence supports. Their small
probability produces an extension of the entire channel into two
outputs. Iteration and a Bell-sharing bound then control every encoder
and decoder. The bound of Lemma~\ref{f:bell} equals Werner's optimal
entanglement fidelity for universal \(1\to k\) qubit cloning
\cite{Werner1998,KeylWerner1999}; we include a short proof.

\subsection{Exact conditional channels and initial supports}

Fix a rate \(\eps\) and abbreviate \(\Lambda_L=\Lambda_{L,\eps}\).
Let this be one full round on the complete physical
register. Dephase the retained classical words at the beginning of
the step. Write \(a\) for their measured initial word and \(\omega\)
for one complete joint history of autonomous metadata and exogenous
innovations, including compiler marks. The argument below also applies
to the four-dimensional appendix: put \(d_{\rm lat}=3\), \(n_{\rm sel}=1\)
for Theorem~\ref{thm:detailed}, and \(d_{\rm lat}=4\), \(n_{\rm sel}=4\)
for its folded part. The latter history includes all selectors and the
profile together. Syndrome measurements and their outcomes stay
inside the quantum CPTP circuit. Conditional on \((a,\omega)\),
the controller trajectory is fixed, and the data evolution is a
CPTP map \(T_{a,\omega}\). All auxiliary hardware is reset before
reading and discarded or reset after the final gatekeeper block.

Draw a conservative dependence graph for this conditional circuit.
A preparation has no incoming quantum edge. An ordinary local CPTP
map has edges from every input factor to every output factor. A clean
bottom-triggered gatekeeper, including its discarded swap register,
has no incoming edge to its retained outputs. Each physical replacement
also removes the incoming edges of the factors it replaces. Include temporary quantum work
until its actual discard. Let \(A_s(a,\omega)\) be the set of initial
physical sites supporting any retained quantum factor that has a path to any retained output after \(s\)
rounds.

\begin{lemma}[Initial-support factorization]\label{f:factorization}
For every \((a,\omega)\), there is a CPTP map \(G_{a,\omega}\) such
that
\begin{equation}\label{f:factor}
 T_{a,\omega}
 =G_{a,\omega}\circ
   \operatorname{Tr}_{A_s(a,\omega)^c}.
\end{equation}
This identity holds on the complete initial quantum algebra and after
tensoring with an arbitrary reference. The fixed conditional final
classical word and all retained hardware outputs are included.
\end{lemma}
\begin{proof}
Propagate an arbitrary output observable backward in the Heisenberg
picture. A CPTP gate's adjoint is unital, so gates outside the
observable's backward support contribute identities. At a
preparation or complete erasure, backward dependence terminates.
By the graph definition the resulting observable is supported on
\(A_s(a,\omega)\). This is true for every output observable, proving
that the conditional channel depends only on the reduced input on
that initial set. Fix a reference density matrix on its complement
and define \(G_{a,\omega}\) by inserting that reference matrix before
\(T_{a,\omega}\). It is CPTP and gives \eqref{f:factor}.

The statement uses the complete cut in every round. In the
concrete schedule all spatial quantum work precedes the final
controlled-swap and discard block. What follows consists of controller
storage, preparation, reset, and noise, so no old quantum state can
be returned from temporary hardware to retained data. Conditional
syndrome processing remains part of a CPTP map throughout this
argument.
\end{proof}

\begin{proposition}[Halting ancestor tail]\label{f:ancestor}
If the computation indexed by \(n\) halts, there are finite positive
constants \(B_n,\kappa_n\) and an integer \(d_n\), uniform over
\(\eps\in I_n\), such that for every measured initial word \(a\)
and initial site \(v\),
\begin{equation}\label{f:ancestor-tail}
 \Pr\{v\in A_s(a,\omega)\}
 \le B_n(1+s)^{d_n}e^{-\kappa_ns}
\end{equation}
on every sufficiently large torus whose depth-\(s\) determining cones
unwrap injectively. The event is determined by innovations in a
spatial cone of radius \(O(s)\) during these \(s\) steps.
\end{proposition}
\begin{proof}
A surviving path crosses a retained gatekeeper in every round.
A bottom gatekeeper that is clean cuts it, so at each crossing either
the ordinary checked proposal is live or there is a direct named
local fault. Proposition~\ref{c:late-live} bounds, simultaneously over all
initial controller words, the possibility that a sufficiently late
live proposal is outside the raw marked field. Only polynomially
many roots of the selector collection occur in a specified depth-\(s\)
quantum cone. Union over these candidates preserves its exponential
bound.

On the complementary event, each crossing of the actual path is
covered by a rare ordinary proposal of its current sheet or a direct
local fault. The path may change sheet at every round. Within
one selector use the south comparison of Lemma~\ref{c:rare}; its
rare birth needs \(256\) fresh positions and its seed groups have
the eight-class bound. Pack complete per-sheet contours along the
actual physical path, reserving each seed test's entire deterministic
two-time support and restarting below its minimum time minus two.

For the above \(n_{\rm sel}\), \(g\) prescribed typed rare groups, and \(k_{\rm dir}\) distinct
direct physical source names, H\"older across the sheets followed by
Cauchy--Schwarz between rare and direct events gives
\[
 (N_{\rm seed}p^{256})^{g/(16n_{\rm sel})}p^{k_{\rm dir}/2}.
\]
The rare power is \(256/(16n_{\rm sel})\ge4\), and the direct power is \(1/2\).
Their positive exponents beat the finite entropy of the graph
containing every allowed physical quantum template, as in
Proposition~\ref{c:path-tail}. Union over polynomially many
reachable endpoints and the preceding marginal exceptions proves
\eqref{f:ancestor-tail}, uniformly in the combined measured word.
For the torus statement first use the literal finite determining
cone of the event. When that cone is injective, its innovations have
the same joint law as the corresponding finite region of an
independent lattice field. Extend outside this cone with independent
innovations to apply the preceding infinite-lattice bound. A separate
extension may be used for each individual upper bound. Dependence
between different torus block events will instead be read directly
from their actual determining cones.
\end{proof}

The constants in Proposition~\ref{f:ancestor} may use a finite
halting motif. Their existence does not supply a total algorithm
that produces a destruction time from every input \(n\).

\subsection{An elementary finite-dependence lemma}

We give the precise domination argument used to turn the ancestor
tail into a channel extension. Independence from nonneighbors means
independence from their entire vector.

\begin{lemma}[Product domination]\label{f:domination}
Let \((X_j)_{j\in J}\) be \(\{0,1\}\)-valued variables on a finite
dependency graph of maximum degree \(\Delta\). Suppose \(X_j\) is
independent of the vector of all its nonneighbors. Let
\(0<q\le1/2\), put \(r=q/2\), and suppose
\[
 \Pr(X_j=1)\le p\le r^{\Delta+1}
 \quad\hbox{for every }j.
\]
Then their joint law is stochastically dominated by independent
Bernoulli\((q)\) variables.
\end{lemma}
\begin{proof}
Add independent variables \(Z_j\) of Bernoulli parameter \(r\),
independent of \(X\), and let \(\tilde X_j=X_j\lor Z_j\). We first prove,
for every \(j\) and every positive-probability pattern \(w\) on a
set \(K\) excluding \(j\), that
\begin{equation}\label{f:conditional}
 \Pr(X_j=1\mid \tilde X_K=w)\le r.
\end{equation}
Induct on \(|K|\). The empty case follows from \(p\le r\). Split
\(K=F\sqcup H\) into nonneighbors and neighbors of \(j\), with
\(|H|\le\Delta\). Vector independence gives
\[
 \Pr(X_j=1,\tilde X_K=w)
 \le p\Pr(\tilde X_F=w_F).
\]
Expose the assignments on \(H\) in any order, after those on \(F\).
A prescribed one has conditional probability at least \(r\), since
its own independent \(Z\)-coin forces it. A prescribed zero has
conditional probability at least \((1-r)^2\): the induction
hypothesis bounds the probability of its \(X\)-variable being one
by \(r\), since the conditioning set at that stage has size less
than \(|K|\), and its own \(Z\)-coin is still independent. Since
\(r\le1/4\), we have \((1-r)^2\ge r\). Hence
\[
 \Pr(\tilde X_K=w)\ge r^{|H|}\Pr(\tilde X_F=w_F).
\]
Dividing gives a ratio at most
\(p r^{-\Delta}\le r\), completing the induction.

The own \(Z_j\)-coin is independent of \((X_j,\tilde X_K)\), so
\[
 \Pr(\tilde X_j=1\mid \tilde X_K)
 \le r+(1-r)r\le q.
\]
Generate the finite vector \(\tilde X\) sequentially from independent
uniform variables, using its conditional probabilities. With the
same uniforms, the tests against \(q\) generate an independent
Bernoulli\((q)\) vector \(Y\) satisfying \(\tilde X\le Y\). Finally sample
\(X\) conditionally on \(\tilde X\) in its original joint law. Then
\(X\le \tilde X\le Y\) and its marginal law is correct. This is the
claimed finite monotone coupling.
\end{proof}

\subsection{Two outputs with the same full-channel marginal}

\begin{proposition}[Two-output extension]\label{f:extension}
For a halting input \(n\), there are finite integers \(s_n,A_n'\),
uniform throughout \(I_n\), such that on every torus
\(L>A_n's_n\) the complete channel \(\Lambda_L^{s_n}\) has a CPTP
extension to two copies of its whole output algebra with both
marginals equal to \(\Lambda_L^{s_n}\).
\end{proposition}
\begin{proof}
Partition the initial torus into coordinate blocks whose sides lie
between \(s\) and \(2s\), which is possible because \(L\) exceeds a
fixed multiple of \(s\). Let \(X_I(a,\omega)\) indicate that block
\(I\) contains an initial quantum factor in \(A_s(a,\omega)\).
Proposition~\ref{f:ancestor} and a union over its \(O(s^{d_{\rm lat}})\) sites
give
\[
 \Pr(X_I=1\mid a)\le
 p_s= B_n(1+s)^{d_n'}e^{-\kappa_ns}\longrightarrow0.
\]
Each indicator is the mask of all quantum factors at the initial physical sites.
The joint innovations determining a block lie in an enlarged block of
radius \(O(s)\). Thus these indicators have a dependency graph of
some degree \(\Delta\) independent of \(s,L,a\). Conditional on
\(a\), nonadjacent determining regions use disjoint underlying
innovations, giving the vector independence in
Lemma~\ref{f:domination}.

Choose \(s=s_n\) large enough that
\(p_s\le(1/4)^{\Delta+1}\), uniformly in the rate window. For a
fixed \(a\), its mask law \(\mathbb P_a\) is dominated by the independent
fair-mask law \(\pi\). Choose a finite monotone coupling and write
\(K_a(X\mid Y)\) for the conditional law of its lower mask given
its upper one. Sample a fair mask \(Y\); both \(Y\) and its
complement \(1-Y\) have law \(\pi\). Sample
\[
 X^{(1)}\sim K_a(\,\cdot\mid Y),\qquad
 X^{(2)}\sim K_a(\,\cdot\mid1-Y).
\]
Their supports are disjoint and each has the exact original mask
law \(\mathbb P_a\). Given these masks, sample two complete joint
innovation histories of every autonomous field and all shared marks with the original conditional law given the
corresponding mask. Each history retains all physical correlations and its exact original marginal
law. Their actual initial ancestor sets are disjoint, because each
lies within its own lower mask.

Route every quantum factor at those disjoint initial physical sites to the two maps
\(G_{a,\omega_1}\) and \(G_{a,\omega_2}\) from
Lemma~\ref{f:factorization}, and trace unused factors. Applying the
product of these maps is CPTP even on a globally entangled input.
For either output separately, averaging its exact conditional
history law and using \eqref{f:factor} recovers the original
\(s_n\)-step channel in branch \(a\).

To obtain a channel on the full initial algebra, first measure the
initial classical register, broadcast its classical outcome \(a\)
to both simulations, and perform the conditional construction above.
This is precisely the dephasing already present in the physical
channel. It therefore preserves its action also when the encoded
input entangles the classical register, data, and reference.
All final controller words and reset hardware outputs belong to
these simulations. This proves the full-algebra extension.
\end{proof}

The extension is a mathematical information bound. Its construction
may use global conditional sampling; the physical processor itself
is unchanged.

\begin{lemma}[Iteration on the entire algebra]\label{f:iteration}
If a channel \(T\) has a two-output extension with equal marginals,
then \(T^j\) has a \(2^j\)-output extension with equal marginals
for every positive integer \(j\).
\end{lemma}
\begin{proof}
Apply the extension separately to every output of the previous
extension. Partial trace over the other outputs commutes with these
trace-preserving maps, so each new marginal is \(T\) composed with
the previous marginal. By induction the depth-\(j\) construction
has \(2^j\) outputs and marginal \(T^j\). Correlations between
previous outputs cause no difficulty for tensor products of CPTP
maps.
\end{proof}

For \(T=\Lambda_L^{s_n}\), this gives \(2^j\) outputs at physical
time \(js_n\). Once the extension exists on a fixed torus, its
iteration uses the full finite algebra; the size requirement remains
\(L>A_n's_n\).

\subsection{Bell sharing and the volume-uniform time}

\begin{lemma}[Bell-sharing bound]\label{f:bell}
Suppose a channel has a \(k\)-output extension with equal marginals.
For every qubit encoder and common qubit decoder, its entanglement
fidelity is at most
\[
 \frac12+\frac{1}{2k}.
\]
\end{lemma}
\begin{proof}
Apply the encoder to the extension's input and the decoder to each
output. Let \(R\) be the reference qubit and \(B_1,\ldots,B_k\)
the decoded qubits. Write \(P_j\) for the normalized Bell projector
on \(RB_j\), tensored with the identity on the other outputs.
Let \(V_j\) insert this Bell state into \(RB_j\); then
\(P_j=V_jV_j^*\). The nonzero spectrum of \(S=\sum_jP_j\) equals
that of its block Gram operator \(G=(V_i^*V_j)_{i,j}\).
The diagonal blocks are identity. Contracting two normalized Bell
tensors shows that for \(i\ne j\), \(V_i^*V_j\) is one half of a
unitary relabeling of the missing-output tensor factors, so its norm
is \(1/2\). For a block vector \(z\),
\begin{align*}
 \langle z,Gz\rangle
 &\le\sum_j\|z_j\|^2+
       \frac12\sum_{i\ne j}\|z_i\|\,\|z_j\|\\
 &\le\frac{k+1}{2}\sum_j\|z_j\|^2.
\end{align*}
Thus \(\|S\|\le(k+1)/2\). The equal decoded marginals have the
same entanglement fidelity \(F\), whence
\(kF=\operatorname{Tr}(\rho S)\le(k+1)/2\).
\end{proof}

\begin{theorem}[Uniform quantum forgetting on a halting window]
\label{f:forgetting}
If the computation indexed by \(n\) halts, then for every \(\eta>0\)
there is a finite \(C_{n,\eta}\), uniform over all torus sides \(L\)
and all \(\eps\in I_n\), such that every qubit encoder and decoder
through \(\Lambda_L^t\) has entanglement fidelity at most
\(1/2+\eta\) for every integer \(t\ge C_{n,\eta}\).
\end{theorem}
\begin{proof}
It suffices to consider \(0<\eta<1/2\). Choose \(j\) with
\(2^{-j-1}\le\eta\). Proposition~\ref{f:extension} and
Lemmas~\ref{f:iteration}--\ref{f:bell} give the bound at time
\(js_n\) on every \(L>A_n's_n\), uniformly in the rate window.

There are finitely many remaining tori. Let \(N_n\) bound their
physical site counts and let \(N\) be the fixed number
of physical qubits per site. At a fixed final elementary noise layer,
simultaneous replacement of every factor has probability at least
\[
 p_n=(\inf I_n)^{R_n}>0,\qquad
 R_n=\begin{cases}N_n,&\mathrm W,\\
 N N_n,&\mathrm Q.\end{cases}
\]
It replaces every physical factor, including controller and work,
and disconnects the output from the reference. Conditional on at
least one such event, the decoded entanglement fidelity is exactly \(1/4\):
the reference marginal is maximally mixed, and the output is
independent of it. The complementary histories have entanglement fidelity at
most one. Independence of the flags across rounds gives
\[
 F_L(t)\le\frac14+\frac34(1-p_n)^t.
\]
Choose \(t_n\) where this is at most \(1/2+\eta\), and let
\(C_{n,\eta}=\max\{js_n,t_n\}\).

An upper bound valid for every decoder persists under later
evolution: compose any proposed later decoder with that additional
evolution to obtain a decoder at the earlier time. Thus the bound
holds at every time in the statement.
\end{proof}
